%%%%%%%%%%  scan idx 649  |  folio 636  %%%%%%%%%%
\SGAPrintPage{636}%
following easy lemma, hence for \(X_{S'}/S'\) by 3.5 and 3.4 (ii), and finally for \(X/S\) by 3.4 (i).

\sgaStableTarget{sga6.e13.lem.3-7}\textbf{Lemma 3.7.} \emph{Let \(X=\coprod X_i\) be of finite type over noetherian \(S\), and \(\mathcal F\) a family of classes of coherent sheaves on the fibers of \(X/S\). For \(\mathcal F\) to be bounded, it is necessary and sufficient that for every \(i\) the induced family \(\mathcal F_i\) on the fibers of \(X_i/S\) be bounded.}

In the proof of 3.6, we reduced to the case where \(X\) is projective and integral over \(S\)\,; endow \(X/S\) with a very ample sheaf \({O}_X(1)\). Let \(\mathcal L\) be a bounded family of classes of invertible sheaves on the fibers of \(X/S\), and \(L\) an invertible sheaf on an \(X_T\), with \(T\) of finite type over \(S\), which bounds \(\mathcal L\). After decomposing \(T\) into pieces, one may suppose \(L\) flat over \(T\), and after replacing \(T\) by each of its connected components, one may suppose \(T\) connected. Then the polynomial in \(q\) and \(m\), \(P(q,m)=\chi(L_t^{\otimes q}(m))\), does not depend on \(t\in T\) (EGA III 7.9.4).

Let \(\mathcal L_n=\Phi_n^{-1}(\mathcal L)\): \(\mathcal L_n\) is the family of classes of invertible sheaves \(M_K\) such that the class of \(M_K^{\otimes n}\) belongs to \(\mathcal L\). Let \(\chi(M_K^{\otimes p}(m))=\sum_{i=0}^{r}a_i(p)\binom{m+i}{i}\), where \(a_i(p)\) is a polynomial in \(p\) of degree \(\le r-i\). For variable \(q,m\), we have \(\chi(M_K^{\otimes qn}(m))=P(q,m)\), so each \(a_i(qn)\) is independent of the chosen \(M_K\)\,; consequently, each polynomial \(a_i(p)\) is also independent of \(M_K\). Therefore the family \(\mathcal L_n\) has a single Hilbert polynomial \(\chi(M_K(m))\), hence is bounded by 2.11.

\sgaStableTarget{sga6.e13.thm.3-8}\textbf{Theorem 3.8.} \emph{Let \(S\) be a noetherian scheme, \(Y\) a projective \(S\)-scheme endowed with a relatively ample invertible sheaf \({O}_Y(1)\), \(X\) the zero scheme of a section \(\varphi\) of \({O}_Y(1)\), and \(f:X\to Y\) the inclusion morphism. We}
%%%%%%%%%%  scan idx 650  |  folio 637  %%%%%%%%%%
\SGAPrintPage{637}%
\emph{suppose that for every point \(s\in S\) and every irreducible component \(Y_s^i\) of \(Y_s\), the irreducible components of \(X_s\cap Y_s^i=V(\varphi|Y_s^i)\) have dimension \(\ge2\), (which is the case if the irreducible components of the fibers of \(Y\) have dimension \(\ge3\)). Then the induced morphism between the Picard functors}
\[
f^*:\underline{\Pic}_{Y/S}\longrightarrow\underline{\Pic}_{X/S}
\]
\emph{is of finite type.}

\emph{Moreover, if \(S\) is integral, there exists a nonempty open subset \(U\) of \(S\) such that the Picard functors restricted to \(U\) are representable and the induced morphism}
\[
\underline{\Pic}_{Y/S}|U\longrightarrow\underline{\Pic}_{X/S}|U
\]
\emph{is quasi-affine (of finite type).}

Let us first prove the first assertion. As in the proof of 3.6, one sees by noetherian induction that one may replace \(S\) by any nonempty open subset, and one easily reduces to the case where \(S\) is integral and \(Y\) is integral over \(S\), always replacing \(X\) by \(X\times_YY'\) when \(Y\) is replaced by \(Y'\)\,; this done, one reduces in the same way to the case where \(Y\) is normal over \(S\), using the following lemma in place of XII 1.3.

\sgaStableTarget{sga6.e13.lem.3-9}\textbf{Lemma 3.9.} \emph{Let \(S\) be an integral noetherian scheme and \(X\) an \(S\)-scheme, integral over \(S\). Then there exist a nonempty open subset \(U\) of \(S\), an integral scheme \(S'\) finite and faithfully flat over \(U\), an \(S'\)-scheme \(X'\) normal and integral over \(S'\), and a finite surjective \(S'\)-morphism \(X'\to X_{S'}\).}

%%%%%%%%%%  scan idx 651  |  folio 638  %%%%%%%%%%
\SGAPrintPage{638}%
Indeed, let \(\eta\) be the generic point of \(S\) and \(K\) the algebraic closure of \(k(\eta)\). It is clear that \(\Spec(K)\) is the limit of the filtered inverse system \((S_\alpha)\) of integral affine schemes endowed with a finite faithfully flat morphism \(S_\alpha\to U_\alpha\), where \(U_\alpha\) is a nonempty open subset of \(S\). The normalization of \(X_K\) therefore comes from an \(S_\alpha\)-scheme \(X'\) for some \(\alpha\), and the finite surjective morphism \(X'_K\to X_K\) from a finite surjective \(S_\alpha\)-morphism \(X'\to X_{S_\alpha}\) (EGA IV 8.8.2, 8.10.5). Finally, since \(X'_K\) is normal and integral, there exists an open subset \(S'\) of \(S_\alpha\) such that \(X'_{S'}\) is normal and integral over \(S'\) (EGA IV 9.9.5).

\sgaStableTarget{sga6.e13.rem.3-10}\textbf{Remark 3.10.} Retaining the notation and conditions of 3.9, suppose that there exists a resolving morphism \(X'_K\to X_K\) (EGA IV 7.9.1), (for example if \(\dim(X_K)\le2\), by Abhyankar). Then a similar argument shows that there exist a finite faithfully flat morphism \(S'\to U\), a scheme \(X'\) smooth over \(S'\) and a proper surjective birational morphism \(X'\to X_{S'}\).

In the proof of the first part of 3.8, we reduced to the case where \(S\) is integral and \(Y\) is normal and integral over \(S\), so \(Y\) is integral. If \(\varphi=0\), then \(X=Y\) and the assertion is trivial. Otherwise, after restricting \(S\), one may suppose that for every \(s\in S\) we have \(\varphi_s\ne0\), so \(\varphi_s\) is \({O}_{Y_s}\)-regular and the dimension \(r\) of \(Y_s\) is \(\ge3\).

Now let \(\mathcal L\) be a bounded family of classes of invertible sheaves on the fibers of \(X/S\). Let \(M_K\) be an invertible sheaf on a geometric fiber \(Y_K\) of \(Y/S\), with Hilbert polynomial \(\chi(M_K(n))=\sum_{i=0}^{r}a_i\binom{n+i}{i}\); then the Hilbert polynomial of \(f^*M_K\) is \(\sum_{i=0}^{r-1}a_{i+1}\binom{n+i}{i}\). As the class of \(M_K\) ranges over \((f^*)^{-1}(\mathcal L)\), that is, the
%%%%%%%%%%  scan idx 652  |  folio 639  %%%%%%%%%%
\SGAPrintPage{639}%
class of \(f^*M_K\) belongs to \(\mathcal L\), the three coefficients \(a_r,a_{r-1},a_{r-2}\) remain bounded by 1.13.

Let \(H={O}_Y(m)\) be a multiple of \({O}_Y(1)\) which is very ample. Then, in the Hilbert polynomials \(\chi(M_K(mn))=\sum_{i=0}^{r}a_i^{(m)}\binom{n+i}{i}\) relative to \(H\) of the \(M_K\) in \((f^*)^{-1}(\mathcal L)\), the three coefficients \(a_r^{(m)},a_{r-1}^{(m)},a_{r-2}^{(m)}\) remain bounded by 2.10. Consequently, \((f^*)^{-1}(\mathcal L)\) is bounded by 2.11\,; this gives the first part of 3.8.

As for the second part of 3.8, the ``generic'' representability of the Picard functors was established in XII 1.2, as was their separatedness. By an argument similar to that of XII 2, using 3.9 and XII 1.1, one reduces to the case of \(Y\) normal over \(S\). Now, for the morphism \(\underline{\Pic}_{Y/S}\to\underline{\Pic}_{X/S}\) to be quasi-affine, it suffices that it be (separated and) quasi-finite (EGA IV 8.12.8), which can be checked on the geometric fibers, and thus follows from the following lemma:

\sgaStableTarget{sga6.e13.lem.3-11}\textbf{Lemma 3.11.} \emph{Let \(k\) be an algebraically closed field, \(Y\) a normal projective \(k\)-scheme endowed with an ample sheaf \({O}_Y(1)\), and \(X\) the zero scheme of a section \(\varphi\) of \({O}_Y(1)\). We suppose that \(Y\) has dimension \(\ge3\) (resp. \(\ge2\)). Then the morphism \(\Pic_{Y/k}\longrightarrow\Pic_{X/k}\) (resp. the morphism \(\Pic^\tau_{Y/k}\to\Pic^\tau_{X/k}\) (see 4.1)) is of finite type and has a kernel \(N\) which is a finite group (in general nonreduced).}

Indeed, one may evidently suppose \(\varphi\ne0\)\,; thus \(\varphi\) is regular, \(Y\) being integral. The morphism under consideration is of finite type by the first part of 3.8 (resp. by 4.7 (iii)). Consequently, the kernel \(N\) is an algebraic group, hence contained in \(\Pic^\tau_{Y/k}\), and is therefore the same
%%%%%%%%%%  scan idx 653  |  folio 640  %%%%%%%%%%
\SGAPrintPage{640}%
in the two cases under consideration. Since \(\Pic^\tau_{Y/k}\) is proper, \(Y\) being normal (FGA 236-2.13), and since \(N\) is a subgroup, necessarily closed (SGA 3 VI\(_B\) 1.9.2), for \(N\) to be finite it suffices that it contain no subgroup of type \(\mu_\ell(=\mathbb Z/\ell\mathbb Z)\) where \((\ell,\operatorname{car}(k))=1\).

Consider Kummer theory, that is, the following commutative diagram, coming from the exact sequence
\[
0\longrightarrow\mu_\ell\longrightarrow\mathbb G_m
\xrightarrow{x\mapsto x^\ell}\mathbb G_m\longrightarrow0
\]
of étale sheaves on \(Y\) (resp. on \(X\)):
\[
\begin{tikzcd}[column sep=1.6em,row sep=1.2em]
& & H^1(Y,\mu_\ell)\arrow[r]\arrow[d]&H^1(Y,\mathbb G_m)\arrow[r,"\ell"]\arrow[d]&H^1(Y,\mathbb G_m)\arrow[d]\\
H^0(X,\mathbb G_m)\arrow[r,"\ell"]&H^0(X,\mathbb G_m)\arrow[r]&H^1(X,\mu_\ell)\arrow[r]&H^1(X,\mathbb G_m)\arrow[r,"\ell"]&H^1(X,\mathbb G_m).
\end{tikzcd}
\]
By Hilbert's ``Theorem 90'', we have \(H^1(Y,\mathbb G_m)=\Pic_{Y/k}(k)\) (resp. \(H^1(X,\mathbb G_m)=\Pic_{X/k}(k)\)). Since \(X\) is proper and integral, we have \(H^0(X,\mathbb G_m)=k^*\)\,; therefore \(H^1(X,\mu_\ell)\to H^1(X,\mathbb G_m)\) is injective.

It therefore suffices to see that \(H^1(Y,\mathbb Z/\ell\mathbb Z)\longrightarrow H^1(X,\mathbb Z/\ell\mathbb Z)\) is injective, which follows from \(\pi_1(X,\xi)\to\pi_1(Y,\xi)\) being surjective, or again (SGA 1 V 6.9) from every connected étale covering \(\widetilde Y\) of \(Y\) with group \(\mathbb Z/\ell\mathbb Z\) having a restriction \(\widetilde X=X\times_Y\widetilde Y\) which is also connected; this is a well-known consequence of Bertini's theorem and Zariski's connectedness theorem (SGA 1 X 2.10).

\sgaStableTarget{sga6.e13.rem.3-12}\textbf{Remark 3.12.} Retain the conditions of 3.11. Then:
\begin{enumerate}[label=(\roman*)]
\item The kernel \(N\) is a finite \uline{unipotent} group (SGA 3 XVII 1.1), since it contains no subgroup of multiplicative type \(\mu_n\) for any \(n\) (SGA 3 XVII 4.6.1 vi), which is proved in the same way, using instead
